\begin{proof}
For (i), diagonalize
$P_f^{1/2}Q_fP_f^{1/2}=U\Sigma_H^2U^*$ and take
$S=P_f^{1/2}U\Sigma_H^{-1/2}$. Direct substitution proves both
identities. Stability of a proposed truncation is an explicit condition,
not a consequence asserted by this proof.

For (ii), subtract the perturbation from the Stein equality. The
operator-norm bound gives the first inequality in
\eqref{I:U:eq:dyadic-verified-margin}. Since $X\succeq I$ and
$I\succeq X/\|X\|$, the positive map $\Phi_{\widehat T}$ contracts
the order-unit norm induced by $X$ by at most the displayed factor.
This bounds its spectral radius. Factoring each cross term through
$X^{1/2}$ proves the sufficient bound.

For (iii), choose dyadic approximations converging entrywise to the
algebraic input. Their $\theta$ tends to zero, so all sufficiently fine
approximations are stable. Form the joint error realization
\[
 B_a=\Udiag(T_a,\widehat T_a),\qquad
 \beta=(b^\top,\widehat b^\top)^\top,\qquad
 \gamma=(c^\top,-\widehat c^\top)^\top.
\]
Its Gramians $P_E,Q_E$ are obtained from the two joint Stein equations.
The block diagonal matrix $\Udiag(X,X)$ supplies a strict Lyapunov
inequality in a neighborhood of the unperturbed tuple. Thus the joint
Stein inverses and $\operatorname{tr}(P_EQ_E)$ depend continuously on
the perturbation there. At zero perturbation the error series vanishes,
so this trace is zero. Consequently it eventually becomes strictly
smaller than $\varepsilon^2$. Algebraic arithmetic can certify this
strict inequality and the Lyapunov margin, and
\[
 \|H_g-H_{\widehat g}\|^2
 \le\|H_g-H_{\widehat g}\|_{\rm HS}^2
 =\operatorname{tr}(P_EQ_E)<\varepsilon^2.
\]
Enumerating increasing precisions with these tests therefore terminates.
\end{proof}
